\chapter{Introduction to Modernization of Indian Tradition}

\section{Modernization as a Theory}

The next topic --- \textbf{``Modernization of Indian Tradition''} (Yogendra Singh) --- is the topic around which the sociology of India revolves. To understand it, one must first understand modernization. Modernization is a \textbf{process of social change} --- ongoing, and cannot be captured in a fixed frame. Its roots lie in \textbf{Weber, Parsons and Marx}; its key proponents are \textbf{Rostow} (traditional society to mass consumption) and \textbf{Daniel Lerner}.

\begin{CriticalPoint}
Our common-sense idea of modernization (highways, metro, English medium, IAS/IPS, fast food, live-in relationships) is actually the \textbf{Western idea} of modernization --- a form of \textbf{cultural hegemony}. Marx's modernization means simply \textbf{progress}, which can even mean a classless society with equal access and no private property.
\end{CriticalPoint}

\begin{QuickRevision}
\begin{itemize}
\item Modernization of Indian Tradition = the core topic (Yogendra Singh).
\item Roots: Weber, Parsons, Marx; proponents Rostow and Lerner.
\item Common-sense modernization = Western idea (cultural hegemony).
\end{itemize}
\end{QuickRevision}

\section{Yogendra Singh \& the Integrated Approach}

In the first chapter of \emph{Modernization of Indian Tradition}, Yogendra Singh reviews the European theories of modernization (Rostow, Lerner) and the theories of social change --- including those originating in India (M. N. Srinivas's Sanskritization and Westernization).

\emph{Modernization of Indian Tradition} is described as a \textbf{confluence (sangam)} where rivers meet --- Yogendra Singh reviewed the Western modernization theories (Lerner, Rostow, Parsons) on one side and the Indian theories of social change (Srinivas's Sanskritization and Westernization) on the other, and brought them together.

\begin{KeyConcept}
Yogendra Singh's most important contribution: an \textbf{Indian approach} to studying Indian society that is \textbf{not anti-Western} --- an \textbf{integrated approach} combining Western and Indian theories. India is the making of the West, so it cannot be understood by Indian theory alone.
\end{KeyConcept}

\begin{ExampleBox}
To say ``India is modernizing'' is absurd --- India is not one; it consists of diverse little traditions, and each pocket experiences modernization differently. E.g.\ the \textbf{Panchayati Raj Institution} kept the traditional panchayat but imposed modern constitutional/democratic principles (free and fair elections) on top --- a ``sandwich'' of tradition and modernity. In the West traditions were lost; in India traditions were never lost --- they adjusted to modern ideas (structural adjustment).
\end{ExampleBox}

\begin{DataFact}
Yogendra Singh died the previous year; Dipankar Gupta called his death the \textbf{``death of an era in Indian sociology''}.
\end{DataFact}

\begin{QuickRevision}
\begin{itemize}
\item Yogendra Singh: integrated approach (West + India), not anti-Western.
\item India is diverse; traditions adjust to modernity (PRI = sandwich; structural adjustment).
\item Yogendra Singh's death = death of an era (Dipankar Gupta).
\end{itemize}
\end{QuickRevision}

\section{Folk-Urban Continuum \& Little/Great Tradition}

Yogendra Singh also reviewed the American anthropologists of the civilizational approach:
\begin{itemize}
\item \textbf{Robert Redfield} --- the \textbf{folk-urban continuum}.
\item \textbf{Milton Singer} --- the \textbf{little tradition / great tradition} idea (roots traceable to Redfield).
\item \textbf{McKim Marriott} --- the \textbf{interaction} between little and great tradition, for which he coined \textbf{universalization} and \textbf{parochialization}.
\end{itemize}

\begin{KeyConcept}
\textbf{Great tradition} = the national/bigger tradition (great only \textbf{quantitatively}, not qualitatively); \textbf{little tradition} = the local/folk tradition. India has many little traditions --- there is no single Indian tradition.
\end{KeyConcept}

\begin{QuickRevision}
\begin{itemize}
\item Redfield: folk-urban continuum; Singer: little/great tradition; Marriott: universalization \& parochialization.
\item Great tradition = national (quantitative); little = local/folk; India = many little traditions.
\end{itemize}
\end{QuickRevision}
